% ============================================================================
\section{Determinación de los objetivos de la investigación}

\subsection{Objetivo general}

Evaluar de manera rigurosa y transparente si un modelo de computación de reservorio cuántico predice mejor las revisiones de la primera cifra del PBI peruano que los modelos estadísticos y de aprendizaje automático clásicos, usando pronósticos mensuales en tiempo real y simulaciones calibradas.

\subsection{Objetivos específicos}

\begin{enumerate}
	\item Auditar y ordenar el archivo histórico de publicaciones del PBI peruano, identificando retrasos reales de difusión, cambios metodológicos y cobertura sectorial.
	\item Describir estadísticamente las revisiones para saber si presentan sesgo hacia arriba o abajo, dispersión cambiante, ciclos repetitivos o efectos por sector económico.
	\item Entrenar y medir un grupo variado de modelos tradicionales y de aprendizaje automático clásicos como referencia justa de comparación.
	\item Implementar el modelo de reservorio cuántico y compararlo en igualdad de condiciones frente a reservorios clásicos equivalentes y referentes lineales.
	\item Simular doce escenarios artificiales inspirados en el Perú para descubrir en qué tipo de datos podría destacar cada modelo.
	\item Aplicar pruebas estadísticas conjuntas para distinguir mejoras reales de simples coincidencias por azar o por épocas atípicas.
	\item Ejecutar una valoración económica práctica únicamente si el modelo cuántico demuestra una superioridad clara y consistente frente a las alternativas clásicas.
\end{enumerate}
